\chapter{Results Summary}
\label{ch:results}

In the previous chapter we have presented the data analysis that was performed on the 2024 data to extract the hyperfine splitting of the antihydrogen atom. The analysis has been developed with blinded results, and the analysis results have been revealed in a confidential collaboration meeting in the early months of 2025.

The hyperfine splitting has been measured for the two different datasets, and the final results are:

\begin{equation*}
\nu^{high field}_{\overline{H}} = 1420405.8 \pm 1.7 \pm 7.0 \, kHz
\end{equation*}

for the dataset taken at higher magnetic field values (\SI{1.07}{\tesla}), while 

\begin{equation*}
\nu^{high field}_{\overline{H}} = 1420403.9 \pm 1.5 \pm 8.8 \, kHz
\end{equation*}

for the dataset taken at lower magnetic field values (\SI{1.03}{\tesla}). The results have been statistically combined, with a weighted mean, also computing the combined systematic error of the two series (as explained in Sec. \ref{sec:ErrorBudgetSec}). This result represents the most precise determination of the hyperfine splitting of the ground state of the antihydrogen atom. This experimental result has been published in the journal Nature \cite{ALPHA:2026lbu}. The combined result is

\begin{equation*}
\nu_{\overline{H}} = 1420404.8 \pm 1.1 \pm 5.6 \, kHz
\end{equation*} 

\section{Table of Results}

Tab \ref{tab:resultsBdrift} contains the numerical values of the residuals from the linear common fit extraction, as discussed in Sec. \ref{sec:CommonFitExtraction}. The table contains the magnetic field decay $m$ determined with the two methods ("common" fit and "split" method), with the first one used in the nominal approach and the latter for systematic evaluation.

\begin{table}[h]
\centering
\resizebox{\columnwidth}{!}{
\begin{tabular}{l|cc|cc}
\hline
\hline
\textbf{Residuals (kHz)} & \textbf{Series 7} (\cb)  &  \textbf{Series 7} (\da) & \textbf{Series 8} (\cb) &  \textbf{Series 8} (\da) \\ 
\hline 
$\Delta t$ (s) & \multicolumn{2}{c|}{812.4} & \multicolumn{2}{c}{812.4} \\
\hline \hline 
$m_{common}$ (kHz/h) & \multicolumn{2}{c|}{$-75.64 \pm 0.05$} & \multicolumn{2}{c}{$-72.82 \pm 0.04$} \\
\hline 
r0 & $ -16.4 \pm 1.5 $ & $ -15.8 \pm 1.5 $ & $ 6.8 \pm 1.2 $ & $ 0.9 \pm 1.2 $ \\
r1 & $ -0.3 \pm 1.5 $ & $ -0.9 \pm 1.5 $ & $ 6.0 \pm 1.2 $ & $ 3.3 \pm 1.2 $ \\
r2 & $ 9.9 \pm 1.5 $ & $ 6.8 \pm 1.5 $ & $ -2.4 \pm 1.2 $ & $ -5.1 \pm 1.2 $ \\
r3 & $ 10.3 \pm 1.5 $ & $ 3.3 \pm 1.5 $ & $ -11.6 \pm 1.2 $ & $ -1.5 \pm 1.2 $ \\
r4 & $ 3.9 \pm 1.5 $ & $ 8.9 \pm 1.5 $ & $ -9.7 \pm 1.2 $ & $ -8.3 \pm 1.2 $ \\
r5 & $ -0.9 \pm 1.5 $ & $ -2.2 \pm 1.5 $ & $ -0.3 \pm 1.2 $ & $ -4.1 \pm 1.2 $ \\
r6 & $ -2.0 \pm 1.5 $ & $ -7.5 \pm 1.5 $ & $ 3.1 \pm 1.2 $ & $ 9.3 \pm 1.2 $ \\
r7 & $ -0.0 \pm 1.5 $ & $ -10.7 \pm 1.5 $ & $ 5.7 \pm 1.2 $ & $ 2.9 \pm 1.2 $ \\
\hline \hline
$m_{split}$ (kHz/h) & $-75.3 \pm 0.1$  & $-75.9 \pm 0.1$ & $-73.0 \pm 0.1$ & $-72.6 \pm 0.1$\\
%$m_{da}$ (kHz/h) & &  & $-72.6 \pm 0.1$ & $-72.6 \pm 0.1$ \\
%$m_{avg}$ (kHz/h) & $-75.6 \pm 0.1$   & $-75.6 \pm 0.1$ & $-72.8 \pm 0.1$ & $-72.8 \pm 0.1$ \\
\hline
\hline
\end{tabular}
}
\caption{$B$-drift correction parameters and residuals from the common linear fit. The $\Delta t (s)$ value is the time interval between the \fcb and \fda measurement.
\label{tab:resultsBdrift}}
\end{table}


\subsection{Uncertainty Table}

The Tab. \ref{tab:syst} summarizes all the main uncertainties, which have been determined by selecting the greatest contribution from every block in the error budget table \ref{sec:ErrorBudgetSec}. The total systematic uncertainty has been determined as the quadrature sum of each individual contribution. The combination of the two analyzed dataset has been obtained from the sum of the individual shift of each row in \ref{sec:ErrorBudgetSec}, and the total systematic uncertainty of \SI{5.6}{\kilo \hertz} is the resulting from the quadrature sum of each individual contribution. This procedure has been followed to consider possible correlation between the uncertainty of each dataset. The final uncertainty can also be estimated considering the error propagation formula of two independent measurement $$\bar{x} = \frac{x/\sigma_x^2 + y/\sigma_y^2}{1/\sigma_x^2 + 1/\sigma_y^2}, \qquad \sigma_{\bar{x}} = \left(\frac{1}{\sigma_x^2} + \frac{1}{\sigma_y^2}\right)^{-1/2}$$

the result considering uncorrelation is $\sigma_{TOT} = 5.5 $ kHz, in agreement with the nominal estimation of $5.6$ kHz.

\begin{table}[!h]
\centering
\resizebox{\columnwidth}{!}{
\begin{tabular}{r|c|c|c}
\textbf{Source}  & \textbf{\SI{1.07}{\tesla} series}  (kHz) &  \textbf{\SI{1.03}{\tesla} series}  (kHz) & \textbf{Combination}\\
\hline \hline
%\textbf{Statistical} & \textbf{-} & \textbf{-} & \textbf{-}\\
%\hline \hline
Line-shape reproducibility                & 2.8 & 7.0 & 2.7\\
\hline
Bare function ({\it choice})                   & 0.5 & 0.7 & 0.6\\
Bare function ({\it parameters})               & 2.5 & 1.0 & 0.9\\
Bare function ({\it uncertainty increase})     & 3.4 & 3.0 & 2.1\\
\hline
Resolution function ({\it choice})             & 0.3 & 0.6 & 1.1\\
Resolution function ({\it frequency marker})   & 3.5 & 3.5 & 3.4\\
\hline
B-drift ({\it fit procedure})                            & 2.3 & 0.5 & 0.7\\
%\hline
B-drift ({\it long time scale linearity})                  & 0.3 & 0.1 & 0.2\\
%\hline
B-drift ({\it short time scale linearity}) & 1.1 & 1.1 & 1.1\\
\hline
Data selection and binning               & 2.0 & 2.0 & 2.0\\
\hline
\hline
\textbf{Total systematic}                & 7.0 & 8.8 & 5.6\\
\end{tabular}
}
\caption{Summary of the uncertainties associated to the measurement of the hyperfine splitting.
\label{tab:syst}}
\end{table}


\section{CPT, outlook} \label{sec:outlookthesis}

This recent result from the \ALPHA collaboration allows a direct comparison with the most precise determination of the ground state hyperfine level on hydrogen for testing CPT invariance in a model-independent test. While the agreement between the antihydrogen and hydrogen ground-state hyperfine splitting constitutes a stringent, model-independent test of CPT symmetry, any measurable discrepancy would directly signal CPT violation, without reference to any specific theoretical framework. Considering the current uncertainty of the measurement, dominated by the \SI{5.6}{\kilo \hertz} systematic contribution, one can directly determine:

\begin{equation}
\frac{\delta \nu_{\overline{H}}}{\nu_{\overline{H}}} = \frac{5.6 \; kHz}{1420404.8 \; kHz} \approx 4 \times 10^{-6} 
\end{equation}

At the current 4 ppm level of precision, the uncertainty on the antihydrogen ground state hyperfine splitting is now well below the corrections arising from the charge and magnetization distribution of the antiproton, analogous to the well-known proton structure (Zemach radius) which for hydrogen causes a finite size shift of about $-32$ ppm \cite{Widmann_2022,PhysRevA.78.022517}. It is possible to disentangle nuclear structure correction with pure QED contributions by using a separate measurement, the hyperfine splitting of the 2S energy level. The difference $D_{21} = 8 \Delta E_{hfs}(2S) - \Delta E_{hfs}(1S)$, which is called Sternheim interval, has the property of largely cancelling out the nuclear structure terms, isolating an independent quantity for precision QED tests \cite{Sternheim:1963zz}. With this most recent result, the Sternheim interval for antihydrogen has been determined with a new precision:

\begin{equation*}
(8 \Delta E_{hfs}(2S) - \Delta E_{hfs}(1S)) / h = 100 \pm 150 \; kHz
\end{equation*}

where the precision is now limited by the measurement of the $\nu_{2S, \overline{H}}$ splitting, which is $\nu_{2S,\overline{H}} = 177563 \pm 18$ kHz \cite{ALPHA:2026lbu}. The latter was first determined by combining laser spectroscopy of the two accessible hyperfine components of the 1S-2S transition with the previous \ALPHA determination of the ground-state splitting; the present result improves it by a factor 26, \cite{Baker:2025hfs,ALPHA:2026lbu}. The current determination on hydrogen, predicted by QED, leads to $48.9592 \pm 0.0068 \; kHz$ \cite{PhysRevLett.130.203001}, and shows good agreement with the new determination on antihydrogen.

Within the Standard-Model Extension (SME) \cite{ColladayKostelecky1997,ColladayKostelecky1998}, however, the observable reported here has no leading-order sensitivity to Lorentz- or CPT-violating coefficients, and this is a direct consequence of the way it is constructed: the transitions exploited drive the atom from a trapped, low-field-seeking state into an untrapped, high-field-seeking one, and the hyperfine constant is obtained by combining two such field-dependent resonances so as to cancel the local trapping field. The sensitivity of hydrogen and antihydrogen transitions to the various SME coefficients, including the ground-state Zeeman-hyperfine transitions studied in this work, was investigated in \cite{Bluhm_1999,Kostelecky:2015nma}. 

At leading order, the spin-dependent SME operators act on the 1S manifold as one-body perturbations, $\delta H = 2X_e S_z + 2X_p I_z$, with $X_w = -b^w_3 + d^w_{30} m_w + H^w_{12}$ for hydrogen and $b^w_3 \to -b^w_3$ for antihydrogen when the states are labelled, as here, by their energy ordering in the field \cite{Bluhm_1999}\footnote{The coefficient substitution quoted below Eq.~(2) of Ref.~\cite{Bluhm_1999} reads instead $a^w_\mu \to -a^w_\mu$, $d^w_{\mu\nu} \to -d^w_{\mu\nu}$, $H^w_{\mu\nu} \to -H^w_{\mu\nu}$, and refers to the $\ket{m_J,m_I}$ labelling. Since the magnetic moment of the positron is parallel to its spin, the energy ordering of the antihydrogen sublevels is reversed with respect to hydrogen; combining the two effects gives the net sign change of $b^w_3$ quoted here, as stated in Ref.~\cite{Bluhm_1999} for the mixed states $\ket{a}$ and $\ket{c}$.}; the spin-independent coefficients ($a^w_0$, $c^w_{00}$) shift all four sublevels equally. Each spin-dependent term is thus formally an additional, fixed-in-space Zeeman field acting on one spin only. Since $f_{da} - f_{cb} = A\hbar/2\pi$ holds for arbitrary Zeeman couplings of the two spins (see the general expressions retaining the antiproton Zeeman term $\omega_n I_z$ at the end of Ch.~\ref{ch:hyperfine}), the SME terms cancel in the difference for the same reason the trapping field does: $f_{da}$ and $f_{cb}$ shift by the same amount, $[X_e(1+\hat\kappa) + X_p(1-\hat\kappa)]/h$ with $\hat\kappa = \cos 2\theta_1$ the spin-mixing parameter of the $\ket{a}$ and $\ket{c}$ states, at any value of the magnetic field. Consistently with this, $f_{da}-f_{cb}$ coincides with the zero-field $\ket{F=1,m_F=0} \leftrightarrow \ket{F=0,m_F=0}$ maser transition, which is shown in \cite{Kostelecky:2015nma} to be insensitive to propagation-type Lorentz violation for operators of any mass dimension, since only transitions with $\Delta m_F \neq 0$ acquire leading-order shifts. The hydrogen--antihydrogen agreement reported here therefore tests the equality of the bound-state magnetic interaction, including the antiproton finite-size contribution discussed above, but does not constrain at leading order any of the SME coefficients considered in \cite{Bluhm_1999,Kostelecky:2015nma}.

The individual resonances do carry SME information. \fcb and \fda are predominantly a positron spin flip: at \SI{1}{\tesla}, $\hat\kappa \simeq 0.999$, so that $f_{cb}$ and $f_{da}$ shift by $\approx 2X_e/h$, with the antiproton admixture suppressed by $1-\hat\kappa \approx 10^{-3}$. A hydrogen--antihydrogen comparison of either resonance at the same field would therefore isolate the CPT-odd coefficient $b^e_3$ of the electron sector, the CPT-even $d^e_{30}$ and $H^e_{12}$ cancelling in the difference. Exploiting this would require an independent absolute calibration of the field at the location of the atoms; with the 1 ppm accuracy of the ECR technique ($\approx$ \SI{28}{\kilo \hertz} at \SI{1}{\tesla}) the reach on $b^e_3$ would be of order \SI{e-19}{\giga \electronvolt}, several orders of magnitude weaker than the bounds already set on this sector by the comparison of the electron and positron $g$-factors in Penning traps \cite{VanDyck:1987ay}. Despite this, a more promising approach can be constituted by the $c \rightarrow d$ transition, which is mainly an antiproton spin flip. Within the SME framework, it can be shown that hydrogen and antihydrogen would exhibit sensitivity to CPT violating couplings. Specifically, as computed in \cite{Bluhm_1999,Kostelecky:2015nma}, the transition frequency are at first order approximation in $\alpha$

\begin{align}
\delta^{H}_{c \rightarrow d} &\approx ( - b_{3}^{p} + d_{30}^{p} m_{p} + H_{12}^{p}) / \pi \\
\delta^{\overline{H}}_{c \rightarrow d} &\approx ( + b_{3}^{p} + d_{30}^{p} m_{p} + H_{12}^{p}) / \pi
\end{align}

The instantaneous difference of the two transition is therefore $\Delta \nu_{c \rightarrow d} = \nu^{H}_{c \rightarrow d} - \nu^{\overline{H}}_{c \rightarrow d} = - 2 b_{3}^{p}/ \pi$. This result is substantial, a future determination of the magnetic-resonance type transition of the $\ket{c} \leftrightarrow \ket{d}$ between two trapped hyperfine sublevels would isolates the antiproton coefficient $b^{p}_{3}$ probing the proton/antiproton sector of the SME, where the tightest existing antimatter constraints come from the BASE measurement of the antiproton magnetic moment \cite{Smorra2017}.  
At the minimal-SME level this transition probes the same coefficient combination $b^p_3$ already constrained, at the \SI{e-24}{\giga\electronvolt} level, by the BASE comparison of the proton and antiproton magnetic moments \cite{Smorra2017}, so that matching that reach would require an absolute frequency resolution of order \SI{0.1}{\hertz} and hence anti-atom temperatures in the \si{\micro\kelvin} range \cite{Bluhm_1999}. Its interest lies rather in the nonminimal sector, where the bound antiproton momentum $|p|\sim\alpha m_e$ selects combinations inaccessible to Penning-trap experiments \cite{Kostelecky:2015nma}, and in a dedicated search for sidereal variations.



% In addition, since the laboratory-frame coefficients rotate with the Earth, a time-averaged comparison is only sensitive to their component along the Earth's rotation axis, and a dedicated sidereal-time analysis would be needed to access the remaining components \cite{Kostelecky:2015nma}. A measurement with genuine discovery potential would instead need to target a less-constrained sector.

% Moreover, its frequency is stationary with respect to the magnetic field at $B \simeq \SI{0.65}{\tesla}$ \cite{Bluhm_1999}, which removes at first order the Zeeman broadening due to the field inhomogeneity that shapes the present lineshapes and is ultimately responsible for the dominant systematic uncertainties. Since both states are trapped, the transition must be detected through a subsequent microwave transfer to an untrapped state, a double-resonance scheme compatible with the present apparatus. Combined with a dedicated search for sidereal variation, this measurement would represent a natural and well-motivated next step toward exploiting antihydrogen hyperfine spectroscopy as a genuine probe of physics beyond the Standard Model, rather than solely as a model-independent CPT test.